% source: J. S. Milne, "Abelian Varieties" (course notes), v2.00, 2008, www.jmilne.org
% pdf_page: 0023
% doc_page: 17 (Chapter I, Section 3 — end of the proof of Theorem 3.1; "Rational maps into abelian varieties.": Theorem 3.2, Lemma 3.3 and the start of its proof)
% retrieved: 2026-07-17
% transcribed_by: claude-fable-5 (vision, rendered at 200dpi via pdftoppm)

% [running head: 3. RATIONAL MAPS INTO ABELIAN VARIETIES, page 17]

\DeclareMathOperator{\Spec}{Spec}       % Spec(k(V)), Spec(O_Z)
\DeclareMathOperator{\Image}{Im}        % Im(O_{G,e}), image
\DeclareMathOperator{\divisor}{div}     % div(f), divisor of a rational function

% [continuation of the proof of Theorem 3.1 begun on the previous page]

$Z \to V$ is a surjective map from a complete curve onto a nonsingular
complete curve that is an isomorphism on open subsets. Such a map must be an
isomorphism (complete nonsingular curves are determined by their function
fields). The restriction of the projection map $V \times W \to W$ to $Z$
($\simeq V$) is the extension of $\varphi$ to $V$ we are seeking.

The general case can be reduced to the one-dimensional case (using schemes).
Let $U$ be the largest subset on which $\varphi$ is defined, and suppose
that $V \smallsetminus U$ has codimension 1. Then there is a prime divisor
$Z$ in $V \smallsetminus U$. Because $V$ is normal, its associated local
ring is a discrete valuation ring $\mathcal{O}_Z$ with field of fractions
$k(V)$. The map $\varphi$ defines a morphism of schemes
$\Spec(k(V)) \to W$, which the valuative criterion of properness
(Hartshorne 1977, II 4.7) shows extends to a morphism
$\Spec(\mathcal{O}_Z) \to W$. This implies that $\varphi$ has a
representative defined on an open subset that meets $Z$ in a nonempty set,
which is a contradiction. \hfill $\square$

\subsection*{Rational maps into abelian varieties.}

\paragraph{Theorem 3.2.}
\textit{A rational map $\alpha \colon V \dashrightarrow A$ from a
nonsingular variety to an abelian variety is defined on the whole of $V$.}

\paragraph{Proof.}
Combine Theorem 3.1 with the next lemma. \hfill $\square$

\paragraph{Lemma 3.3.}
\textit{Let $\varphi \colon V \dashrightarrow G$ be a rational map from a
nonsingular variety to a group variety. Then either $\varphi$ is defined on
all of $V$ or the points where it is not defined form a closed subset of
pure codimension $1$ in $V$ (i.e., a finite union of prime divisors).}

\paragraph{Proof.}
Define a rational map
\[
  \Phi \colon V \times V \dashrightarrow G, \ (x, y) \mapsto
    \varphi(x) \cdot \varphi(y)^{-1}.
\]
More precisely, if $(U, \varphi_U)$ represents $\varphi$, then $\Phi$ is the
rational map represented by
\[
  U \times U \xrightarrow{\;\varphi_U \times \varphi_U\;} G \times G
    \xrightarrow{\;\mathrm{id} \times \mathrm{inv}\;} G
    \xrightarrow{\;m\;} G.
\]
% [sic: printed "G × G --(id×inv)--> G --(m)--> G"; presumably the middle
%  term should be "G × G", i.e. id × inv : G × G → G × G followed by
%  m : G × G → G]
Clearly $\Phi$ is defined at a diagonal point $(x, x)$ if $\varphi$ is
defined at $x$, and then $\Phi(x, x) = e$. Conversely, if $\Phi$ is defined
at $(x, x)$, then it is defined on an open neighbourhood of $(x, x)$; in
particular, there will be an open subset $U$ of $V$ such that $\Phi$ is
defined on $\{x\} \times U$. After possible replacing $U$ by a smaller open
subset (not necessarily containing $x$), $\varphi$ will be defined on $U$.
% [sic: printed "After possible replacing"; presumably "After possibly replacing" is intended]
For $u \in U$, the formula
\[
  \varphi(x) = \Phi(x, u) \cdot \varphi(u)
\]
defines $\varphi$ at $x$. Thus $\varphi$ is defined at $x$ if and only if
$\Phi$ is defined at $(x, x)$. The rational map $\Phi$ defines a map
\[
  \Phi^* \colon \mathcal{O}_{G,e} \to k(V \times V).
\]
Since $\Phi$ sends $(x, x)$ to $e$ if it is defined there, it follows that
$\Phi$ is defined at $(x, x)$ if and only if
\[
  \Image(\mathcal{O}_{G,e}) \subset \mathcal{O}_{V \times V, (x,x)}.
\]
Now $V \times V$ is nonsingular, and so we have a good theory of divisors
(AG, Chapter 12). For a nonzero rational function $f$ on $V \times V$, write
\[
  \divisor(f) = \divisor(f)_0 - \divisor(f)_{\infty},
\]
% [proof of Lemma 3.3 continues on the next page]
